\chapter{Sampling and Study Design}
\label{ch:sampling}

Chapters 1--4 described the data in front of us. This chapter asks where the data came
from. Every later formula (standard errors, confidence intervals, p-values, regression
coefficients) assumes something about how the numbers were collected. If that assumption
is false, the formula still produces a number, and the number is wrong. No amount of
arithmetic repairs a bad sample or a confounded comparison.

Two questions organise the chapter:
\begin{itemize}
  \item \textbf{Can I generalise?} Does what I saw in my sample hold for the larger group I
  care about? That depends on \emph{how units were selected}.
  \item \textbf{Can I claim cause?} Did the treatment produce the difference, or did
  something else? That depends on \emph{how treatments were assigned}.
\end{itemize}

\section{Populations, parameters, and statistics}
\label{sec:sampling-pop}

Suppose you run a fleet of $6$ cache servers and want the average p50 latency (each
server's median request time) across the fleet. Their true p50 latencies (ms) are
\[
  10,\quad 14,\quad 18,\quad 22,\quad 26,\quad 48 .
\]
The fleet average is $(10+14+18+22+26+48)/6 = 138/6 = 23$ ms. If you can measure all six,
you are done. Usually you cannot: the fleet has $60{,}000$ machines, or measuring one is
expensive, or the ``fleet'' is every user who will ever visit your site. So you measure a
few and reason about the rest.

\begin{definition}[Population, sample, parameter, statistic]
The \textbf{population} is the full set of units you want to learn about. A
\textbf{sample} is the subset you actually observe. A \textbf{parameter} is a number
describing the population (fixed, usually unknown). A \textbf{statistic} is a number
computed from the sample (known, but it changes from sample to sample).
\end{definition}

The rule of thumb for notation: \emph{Greek letters for parameters, Latin letters for
statistics}. Hats also mark estimates: $\hat p$ is ``our estimate of $p$''.

\begin{center}
\renewcommand{\arraystretch}{1.25}
\begin{tabular}{@{}lccl@{}}
\toprule
Quantity & Parameter (population) & Statistic (sample) & Example \\
\midrule
Mean & $\mu$ & $\xbar$ & mean latency \\
Standard deviation & $\sigma$ & $s$ & spread of latency \\
Variance & $\sigma^2$ & $s^2$ & \\
Proportion & $p$ & $\hat p$ & conversion rate \\
Regression slope & $\beta$ (or $\beta_1$) & $b$ (or $\hat\beta_1$) & price per m$^2$ \\
Correlation & $\rho$ & $r$ & height vs.\ weight \\
Size & $N$ & $n$ & fleet vs.\ sample size \\
\bottomrule
\end{tabular}
\end{center}

\begin{figure}[H]
\centering
\begin{tikzpicture}[font=\small]
  \draw[thick, sourceblue, fill=sourceblue!5] (0,0) ellipse (2.9cm and 1.75cm);
  \foreach \x/\y in {0.64/0.73, -0.16/-0.76, 0.23/0.22, 1.38/-1.02, 0.61/-1.12,
    -1.51/-0.85, 2.05/-0.6, -0.72/-1.0, -1.03/0.31, -0.84/-0.57, 1.66/-0.06,
    -0.96/-0.06, 1.3/1.03, -1.58/0.77, 0.13/0.83, -2.13/-0.48, 0.58/1.25,
    -0.83/1.27, 0.23/-0.56, -0.23/-0.23, -0.72/0.86, 1.43/0.58, 0.78/0.24,
    1.92/0.57, 2.27/-0.07, 1.0/0.66, 0.86/-0.24, -0.31/0.45, 0.56/-0.45,
    -2.04/0.2, -0.44/-0.54, -1.21/-1.04, 1.35/-0.6, -0.23/1.25, -1.06/0.99,
    -0.11/0.15, 1.78/-0.84, 0.52/-0.03, 0.09/-1.1, -1.38/-0.39, -0.53/-1.32,
    -1.67/0.03, -0.51/0.09, 2.19/0.29, 1.31/0.17, -2.09/0.65}
    \fill[gray!60] (\x,\y) circle (0.06);
  \foreach \x/\y in {1.38/-1.02, -0.96/-0.06, -0.23/-0.23, -0.31/0.45, -1.58/0.77,
    1.92/0.57}
    \fill[quotientorange] (\x,\y) circle (0.09);
  \node[sourceblue, anchor=south] at (0,1.8) {\textbf{Population} ($N$ units)};
  \node[sourceblue] at (0,-2.1) {parameter: $\mu$, $\sigma$, $p$};
  \draw[thick, quotientorange, fill=quotientorange!6] (7.3,0) circle (1.2cm);
  \foreach \x/\y in {-0.5/0.4, 0.3/0.55, -0.6/-0.3, 0.1/-0.1, 0.55/-0.5, -0.1/-0.65}
    \fill[quotientorange] ($(7.3,0)+(\x,\y)$) circle (0.09);
  \node[quotientorange, anchor=south] at (7.3,1.25) {\textbf{Sample} ($n$ units)};
  \node[quotientorange] at (7.3,-1.55) {statistic: $\xbar$, $s$, $\hat p$};
  \draw[-{Stealth}, thick] (3.0,0.55) to[bend left=18] node[above]{sampling} (6.05,0.55);
  \draw[-{Stealth}, thick, targetgreen] (6.05,-0.55) to[bend left=18]
    node[below]{inference} (3.0,-0.55);
\end{tikzpicture}
\caption{Sampling goes from population to sample; inference goes back. The statistic
is computed; the parameter is what we actually want. How the arrow labelled
``sampling'' was carried out decides whether the return arrow is justified.}
\label{fig:sampling-pop-sample}
\end{figure}

\begin{trapbox}
The ``population'' is defined by the question, not by the data you happen to have. If
your logs contain only users who logged in this week, the population for a question
about ``all customers'' is larger than your logs. Write down the target population
before you look at the sample.
\end{trapbox}

\section{Simple random sampling}
\label{sec:sampling-srs}

\begin{definition}[Simple random sample]
A \textbf{simple random sample} (SRS) of size $n$ from a population of size $N$ is a
sample chosen so that every subset of $n$ units has the same chance of being the one
selected.
\end{definition}

Back to the six servers. Take an SRS of size $n=2$. There are $\binom{6}{2}=15$ possible
pairs, each with probability $1/15$. Here are all $15$ sample means:
\[
  12,\ 14,\ 16,\ 16,\ 18,\ 18,\ 20,\ 20,\ 22,\ 24,\ 29,\ 31,\ 33,\ 35,\ 37 .
\]
(For example, the pair $\{10,14\}$ gives $12$, and $\{26,48\}$ gives $37$.) Their average
is $345/15 = 23$, exactly the population mean $\mu$. Individual samples miss, some by a
lot, but the misses in each direction cancel on average: the high misses (up to $+14$)
are fewer but larger than the low ones (down to $-11$).

\begin{figure}[H]
\centering
\begin{tikzpicture}
\begin{axis}[
  width=12cm, height=4.6cm, axis lines=left,
  xmin=8, xmax=40, ymin=0, ymax=2.9, ytick=\empty,
  xlabel={sample mean $\xbar$ (ms)}, xtick={10,15,20,25,30,35,40},
  clip=false,
]
\addplot[only marks, mark=*, mark size=3pt, color=distblue] coordinates {
  (12,0.3) (14,0.3) (16,0.3) (16,0.7) (18,0.3) (18,0.7) (20,0.3) (20,0.7)
  (22,0.3) (24,0.3) (29,0.3) (31,0.3) (33,0.3) (35,0.3) (37,0.3)};
\addplot[only marks, mark=square*, mark size=3pt, color=trapred] coordinates {(20,1.5)};
\draw[densityred, thick, dashed] (axis cs:23,0) -- (axis cs:23,2.3);
\node[densityred, anchor=south] at (axis cs:23,2.3) {$\mu = 23$ = average of all 15};
\node[trapred, anchor=south east] at (axis cs:22.6,1.75) {``typical'' pair $\{18,22\}$};
\end{axis}
\end{tikzpicture}
\caption{All $15$ equally likely SRS means for $n=2$ from the fleet
$\{10,14,18,22,26,48\}$. They balance at the true mean $23$. A hand-picked
``representative'' pair that avoids the odd server lands at $20$, and would land there
every time.}
\label{fig:sampling-srs-means}
\end{figure}

\begin{definition}[Unbiased]
A statistic is \textbf{unbiased} for a parameter if its average over all possible samples
equals the parameter. Its \textbf{bias} is
$\Bias = \E[\text{statistic}] - \text{parameter}$, where $\E[\cdot]$ (the expected
value) means exactly that average over all possible samples, each weighted by its
probability.
\end{definition}

\begin{recoverybox}
Pause and think: replace the $48$ ms server by one at $100$ ms. Is the average of the $15$
sample means still equal to the new $\mu$? (Yes. Each server still appears in exactly
$5$ of the $15$ pairs, so the grand average is still the plain average of the six servers,
now $190/6 \approx 31.7$ ms.)
\end{recoverybox}

Under simple random sampling, $\xbar$ is unbiased for $\mu$. The calculation above is
the whole proof in miniature: each unit appears in the same number of samples, so each
unit gets equal weight in the grand average. Chapter 6 studies the \emph{spread} of the
$15$ values (the standard error); here the point is the \emph{centre}.

\subsection{Why randomness beats judgement}

The intuitive alternative to random selection is to pick ``representative'' units by
hand. An engineer asked for two typical servers would likely choose $18$ and $22$ and
skip the odd one at $48$. That gives $\xbar = 20$, and it gives $20$ \emph{every time}
the engineer is asked. The error is $-3$ ms, and it never averages out.

\begin{whybox}
Human judgement has preferences: it avoids outliers, favours units that are convenient,
familiar, or look ``normal''. Those preferences are correlated with the thing being
measured, so they produce a bias that repeats on every sample. A random mechanism has no
preferences. It can still produce a bad sample, but bad in an unpredictable direction,
and the size of that chance error is something probability theory can quantify
(Chapter 6). Randomness converts an unknown systematic error into a known random one.
\end{whybox}

In practice an SRS needs a \textbf{sampling frame}: a list of every unit, from which a
random number generator picks. For servers, users, or transactions this list usually
exists in a database.

\begin{lstlisting}[language=Python]
import numpy as np
rng = np.random.default_rng(seed=42)       # fixed seed: reproducible
user_ids = np.arange(1_000_000)            # the sampling frame
srs = rng.choice(user_ids, size=500, replace=False)
# SQL equivalent (approx.): SELECT * FROM users ORDER BY random() LIMIT 500;
\end{lstlisting}

\section{Other sampling designs}
\label{sec:sampling-designs}

Simple random sampling is the reference design. Three variations trade simplicity for
precision or cost; one shortcut abandons randomness altogether.

\subsection{Stratified sampling}

A product has $1000$ users: $800$ on the free tier and $200$ paid. Free users average
$3$ sessions per week (SD $2$); paid users average $15$ (SD $5$). The population mean is
\[
  \mu = \frac{800\cdot 3 + 200 \cdot 15}{1000} = \frac{2400+3000}{1000} = 5.4 .
\]
An SRS of $100$ users might, by bad luck, contain $10$ paid users or $30$. Since paid
users are so different, that luck swings $\xbar$ a lot. Instead, split the population
into the two tiers (\textbf{strata}) and sample each separately: $80$ free and $20$ paid,
in proportion to their sizes. Then combine:
\[
  \xbar_{\text{strat}} = 0.8\,\xbar_{\text{free}} + 0.2\,\xbar_{\text{paid}} .
\]

\begin{definition}[Stratified random sample]
Divide the population into non-overlapping groups (\textbf{strata}) that differ from each
other. Take an SRS within each stratum. Combine stratum means with weights
$W_h = N_h/N$ equal to each stratum's share of the population.
\end{definition}

\begin{figure}[H]
\centering
\begin{tikzpicture}[font=\small]
  \draw[thick, sourceblue, fill=sourceblue!4] (0,0) rectangle (8.4,2.4);
  \node[sourceblue, anchor=south west] at (0,2.4) {Free tier: $N_1=800$};
  \foreach \i in {0,...,15} \foreach \j in {0,...,4}
    \fill[gray!55] ({0.4+0.5*\i},{0.3+0.45*\j}) circle (0.07);
  \foreach \i/\j in {1/0, 4/3, 6/1, 9/4, 11/2, 13/0, 2/2, 7/3, 14/3, 10/0, 3/4, 12/4,
    0/1, 5/0, 8/2, 15/1}
    \fill[quotientorange] ({0.4+0.5*\i},{0.3+0.45*\j}) circle (0.11);
  \draw[thick, structurepurple, fill=structurepurple!5] (9.0,0) rectangle (11.4,2.4);
  \node[structurepurple, anchor=south west] at (9.0,2.4) {Paid: $N_2=200$};
  \foreach \i in {0,...,3} \foreach \j in {0,...,4}
    \fill[gray!55] ({9.4+0.5*\i},{0.3+0.45*\j}) circle (0.07);
  \foreach \i/\j in {1/1, 3/3, 0/4, 2/0}
    \fill[quotientorange] ({9.4+0.5*\i},{0.3+0.45*\j}) circle (0.11);
  \node[anchor=north] at (4.2,-0.15) {SRS of $n_1=80$ (20\%)};
  \node[anchor=north] at (10.2,-0.15) {SRS of $n_2=20$ (20\%)};
\end{tikzpicture}
\caption{Proportional stratified sampling: the same $20\%$ fraction is drawn at random
inside each stratum (orange dots). The paid stratum is guaranteed exactly its fair share
of the sample, so luck in the mix of tiers is removed.}
\label{fig:sampling-stratified}
\end{figure}

How much does this help? Chapter 6 develops standard errors properly; here is the
preview, ignoring the small correction for sampling without replacement. For an SRS the
standard error of $\xbar$ is $\sigma/\sqrt n$, where $\sigma^2$ is the variance of the
whole population. That variance has two parts:
\begin{align*}
  \text{within tiers:}\quad & 0.8\cdot 2^2 + 0.2\cdot 5^2 = 3.2 + 5.0 = 8.2,\\
  \text{between tiers:}\quad & 0.8\,(3-5.4)^2 + 0.2\,(15-5.4)^2 = 4.608 + 18.432 = 23.04,\\
  \text{total:}\quad & \sigma^2 = 8.2 + 23.04 = 31.24 .
\end{align*}
So $\SE_{\text{SRS}} = \sqrt{31.24/100} \approx 0.559$. Stratification removes the
between-tier part entirely, because the tier mix is fixed by design:
\[
  \SE_{\text{strat}} = \sqrt{0.8^2\cdot\frac{2^2}{80} + 0.2^2\cdot\frac{5^2}{20}}
  = \sqrt{0.032 + 0.05} = \sqrt{0.082} \approx 0.286 .
\]
The same $100$ users give roughly half the standard error. To get that from an SRS you
would need about $31.24/0.082 \approx 381$ users.

\begin{recoverybox}
Pause and think: if free and paid users had the same mean ($5.4$ each) but the same SDs
as before, would stratifying still help? (Barely. The between-tier part would be $0$, so
$\sigma^2$ would equal the within-tier $8.2$ and SRS would already achieve
$\SE \approx 0.286$.)
\end{recoverybox}

\begin{intuitionbox}
Stratify on variables that (a) you know for every unit in advance and (b) strongly
predict the outcome. Platform, country, plan tier, and device class are typical. The
more different the strata are from each other, the larger the gain. Stratifying on
something unrelated to the outcome costs nothing but gains nothing.
\end{intuitionbox}

\subsection{Cluster sampling}

To survey on-call engineers about tooling, you could build a list of all $3000$
engineers and pick $120$ at random, which would mean interviewing in $40$ offices.
Instead, pick $4$ teams of $30$ at random and interview everyone on them. The teams are
\textbf{clusters}.

\begin{definition}[Cluster sample]
Divide the population into groups (\textbf{clusters}), take an SRS of clusters, and
measure every unit inside the chosen clusters (or a random subsample of them).
\end{definition}

\begin{figure}[H]
\centering
\begin{tikzpicture}[font=\small]
  \foreach \c in {0,...,11}{
    \pgfmathsetmacro{\cx}{mod(\c,6)*2.0}
    \pgfmathsetmacro{\cy}{-floor(\c/6)*1.7}
    \draw[gray!70, rounded corners] (\cx,\cy) rectangle ++(1.7,1.4);
    \foreach \i in {0,1,2} \foreach \j in {0,1}
      \fill[gray!55] ({\cx+0.35+0.5*\i},{\cy+0.45+0.5*\j}) circle (0.08);
  }
  \foreach \c in {1,4,8}{
    \pgfmathsetmacro{\cx}{mod(\c,6)*2.0}
    \pgfmathsetmacro{\cy}{-floor(\c/6)*1.7}
    \draw[quotientorange, very thick, rounded corners, fill=quotientorange!12]
      (\cx,\cy) rectangle ++(1.7,1.4);
    \foreach \i in {0,1,2} \foreach \j in {0,1}
      \fill[quotientorange] ({\cx+0.35+0.5*\i},{\cy+0.45+0.5*\j}) circle (0.1);
  }
  \node[anchor=north] at (5.85,-1.85) {12 clusters (teams, data centres, schools);
    3 chosen at random; everyone inside is measured};
\end{tikzpicture}
\caption{Cluster sampling picks whole groups. It is cheap when travel or setup per
group dominates the cost, but units in one cluster tend to resemble each other, so $n$
clustered units carry less information than $n$ independent ones.}
\label{fig:sampling-cluster}
\end{figure}

\begin{trapbox}
Strata and clusters are opposites. Strata should be \emph{internally similar and
different from each other}; you sample \emph{from every} stratum. Clusters are ideally
\emph{internally diverse, each a mini-population}; you sample \emph{some} clusters.
Analysing a cluster sample as if it were an SRS understates the standard error, often
badly. Users within the same company, requests from the same session, and pixels from
the same image are all clusters.
\end{trapbox}

\subsection{Systematic sampling}

To audit $10{,}000$ log lines, pick a random start between $1$ and $50$ and then take
every $50$th line. This is a \textbf{systematic sample}. It is simple and spreads the
sample evenly over the list. It behaves like an SRS when the list order is unrelated to
the outcome.

It fails when the list has a period that lines up with the step. Traffic has a daily
cycle. Sampling every $24$th hour hits the same clock time every day.

\begin{figure}[H]
\centering
\begin{tikzpicture}
\begin{axis}[
  width=12cm, height=4.8cm, axis lines=left,
  xmin=0, xmax=72, ymin=0, ymax=190,
  xlabel={hour}, ylabel={requests/s},
  xtick={0,12,24,36,48,60,72}, legend style={font=\small, at={(0.5,1.03)},
  anchor=south, legend columns=2, draw=none,
  /tikz/every even column/.append style={column sep=0.6cm}},
]
\addplot[distblue, thick, domain=0:72, samples=200] {100 + 60*sin(360*(x-6)/24)};
\addlegendentry{true traffic (mean 100)}
\addplot[only marks, mark=*, mark size=3pt, trapred] coordinates {(12,160) (36,160) (60,160)};
\addlegendentry{every 24 h from noon}
\addplot[densityred, dashed, domain=0:72] {100};
\end{axis}
\end{tikzpicture}
\caption{Systematic sampling with a step equal to the cycle length. Every sampled hour is
at the daily peak, so the estimate is $160$ requests/s instead of the true mean $100$.
A random start does not help: whichever start you draw, all points share one phase.}
\label{fig:sampling-systematic}
\end{figure}

\subsection{Convenience sampling}

A \textbf{convenience sample} uses whoever is easiest to reach: the users who answered an
in-app pop-up, colleagues down the hall, the first $1000$ rows of a table sorted by
signup date. There is no random mechanism, so there is no way to compute its error.
Convenience samples are useful for finding bugs and generating hypotheses. They are not
a basis for estimating population numbers.

\begin{center}
\renewcommand{\arraystretch}{1.2}
\begin{tabular}{@{}>{\raggedright\arraybackslash}p{2.4cm}>{\raggedright\arraybackslash}p{4.6cm}>{\raggedright\arraybackslash}p{5.6cm}@{}}
\toprule
Design & How units are chosen & Use when \\
\midrule
SRS & every subset of size $n$ equally likely & you have a full list; default \\
Stratified & SRS inside each stratum & known groups differ on the outcome \\
Cluster & SRS of groups, all units inside & listing or reaching units is costly \\
Systematic & every $k$th unit, random start & list order unrelated to outcome \\
Convenience & whoever is easy & exploration only \\
\bottomrule
\end{tabular}
\end{center}

\section{Bias: errors that do not average out}
\label{sec:sampling-bias}

A statistic can be wrong in two ways. \textbf{Chance error} (variance) changes from sample
to sample and shrinks as $n$ grows. \textbf{Bias} is a systematic offset that comes from
the method and stays the same no matter how many units you collect.

\begin{figure}[H]
\centering
\begin{tikzpicture}[font=\small]
  \foreach \k/\lab/\dx/\dy/\s in {
      0/{low bias\\low variance}/0/0/0.25,
      1/{low bias\\high variance}/0/0/0.75,
      2/{high bias\\low variance}/0.7/0.5/0.25,
      3/{high bias\\high variance}/0.7/0.5/0.75}{
    \begin{scope}[xshift=\k*3.15cm]
      \draw[gray!60] (0,0) circle (1.25); \draw[gray!60] (0,0) circle (0.85);
      \draw[gray!60] (0,0) circle (0.45);
      \fill[densityred] (0,0) circle (0.07);
      \foreach \u/\v in {0.3/0.2, -0.4/0.1, 0.1/-0.5, -0.2/-0.3, 0.5/-0.1,
                         -0.1/0.45, 0.25/-0.2, -0.45/-0.4}
        \fill[distblue] ({\dx+\s*\u*1.6},{\dy+\s*\v*1.6}) circle (0.075);
      \node[align=center, anchor=north] at (0,-1.35) {\lab};
    \end{scope}}
\end{tikzpicture}
\caption{Bias versus variance. Red centre: the parameter. Blue dots: estimates from
repeated samples. Bias is where the cloud is centred; variance is how wide it is. A
larger sample shrinks the cloud but does not move its centre.}
\label{fig:sampling-dartboard}
\end{figure}

Most bias in practice comes from one of four sources.

\subsection{Selection bias}

\textbf{Selection bias} occurs when the way units enter the sample is related to the
outcome. Surveying gym members about exercise habits, measuring app satisfaction among
users who still have the app installed, estimating commute times from people who answered
the survey at the office at 9 a.m.: in each case the selection process prefers certain
values of the thing being measured.

\subsection{Nonresponse bias}

You email a satisfaction survey to $2000$ customers. $800$ reply ($40\%$), and $640$ of
those are satisfied, so $\hat p = 640/800 = 80\%$. What can you say about all $2000$?

The $1200$ non-responders are unknown. In the worst case none of them are satisfied:
$640/2000 = 32\%$. In the best case all are: $(640+1200)/2000 = 92\%$. If non-responders
are satisfied at $50\%$ (plausible: unhappy customers often ignore surveys), the truth is
$(640 + 600)/2000 = 62\%$, which is $18$ points below the headline.

\begin{figure}[H]
\centering
\begin{tikzpicture}[font=\small, x=0.115cm]
  \draw[thick, -{Stealth}] (0,0) -- (103,0);
  \foreach \t in {0,20,40,60,80,100}{
    \draw (\t,0.12) -- (\t,-0.12); \node[below] at (\t,-0.12) {\t\%}; }
  \fill[structurepurple!20] (32,0.25) rectangle (92,0.75);
  \draw[structurepurple, thick] (32,0.25) rectangle (92,0.75);
  \node[structurepurple, anchor=east] at (31,0.5) {possible: 32\%--92\%};
  \draw[quotientorange, very thick] (80,-0.05) -- (80,1.75);
  \node[quotientorange, anchor=south west] at (78,1.75) {responders: 80\%};
  \draw[targetgreen, very thick] (62,-0.05) -- (62,1.75);
  \node[targetgreen, anchor=south east] at (64,1.75) {if non-resp.\ 50\%: 62\%};
\end{tikzpicture}
\caption{With a $40\%$ response rate, the data alone pin the true satisfaction rate only
to a $60$-point-wide range. A large $n$ of responders does not narrow this range; only a
higher response rate does.}
\label{fig:sampling-nonresponse}
\end{figure}

\begin{trapbox}
Reporting $\hat p$ among responders as if it were the population value silently assumes
non-responders look exactly like responders. That is an assumption about data you did
not collect. Report the response rate next to every survey estimate, and treat low rates
($<50\%$) with suspicion.
\end{trapbox}

\subsection{Survivorship bias}

\textbf{Survivorship bias} is selection bias where the filter is ``still around to be
measured'': units that failed, died, or dropped out vanish from the data, and the
survivors look better than the whole group did.

\begin{historybox}
In 1943 the US military was losing bombers over Europe and wanted to add armour. Armour
is heavy, so it could only go where it mattered most. Officers examined bombers
returning from missions and recorded where they had been hit. The holes clustered on the
wings, fuselage, and tail; the engines had few. The natural conclusion was to armour
where the holes were.

Abraham Wald, a Hungarian-born mathematician working at the Statistical Research Group at
Columbia University (a wartime team of statisticians solving military problems), pointed
out the selection: these were the planes that \emph{came back}. Bullets do not prefer
wings to engines. Planes hit in the engines mostly did not return, so they were missing
from the data. The clean areas on survivors marked where hits were fatal, and armour
belonged there. Wald went further and worked out how to estimate the vulnerability of
each area from the survivors alone, an early example of correcting an estimate for data
that are missing for a known reason.
\end{historybox}

\begin{figure}[H]
\centering
\begin{tikzpicture}[scale=0.85, font=\small]
  % wings
  \draw[thick, fill=gray!15] (-0.3,0.6) -- (-4.5,-0.2) -- (-4.5,-0.8) -- (-0.3,-0.4) -- cycle;
  \draw[thick, fill=gray!15] (0.3,0.6) -- (4.5,-0.2) -- (4.5,-0.8) -- (0.3,-0.4) -- cycle;
  % tail
  \draw[thick, fill=gray!15] (-0.2,-2.35) -- (-1.6,-2.85) -- (-1.6,-3.2) -- (-0.2,-3.0) -- cycle;
  \draw[thick, fill=gray!15] (0.2,-2.35) -- (1.6,-2.85) -- (1.6,-3.2) -- (0.2,-3.0) -- cycle;
  % fuselage
  \draw[thick, fill=gray!15] (0,0) ellipse (0.38cm and 3.2cm);
  % engines
  \foreach \ex/\ey in {-1.6/0.2, 1.6/0.2, -3.0/-0.1, 3.0/-0.1}
    \draw[thick, fill=targetgreen!25] (\ex,\ey) ellipse (0.22cm and 0.55cm);
  % cockpit
  \draw[thick, fill=targetgreen!25] (0,2.35) ellipse (0.22cm and 0.4cm);
  % bullet holes
  \foreach \x/\y in {-0.8/0.1, -1.1/-0.2, -2.2/-0.35, -2.4/0.05, -3.6/-0.45,
      -4.1/-0.5, -3.9/-0.3, -2.0/-0.5, -0.7/-0.25,
      0.8/0.2, 1.0/-0.15, 2.3/-0.2, 2.2/-0.5, 3.7/-0.55, 4.0/-0.35, 2.5/0.1,
      0.75/-0.3, 3.6/-0.3,
      -0.1/1.1, 0.12/0.7, -0.15/-0.9, 0.1/-1.3, -0.12/-1.7, 0.15/-2.0,
      0.0/1.5, 0.1/-0.4,
      -0.8/-2.75, -1.2/-2.95, 0.9/-2.8, 1.3/-3.0, -0.5/-2.6}
    \fill[trapred] (\x,\y) circle (0.065);
  \node[anchor=west, align=left, targetgreen!60!black] (eng) at (5.1,1.5)
    {engines, cockpit:\\ almost no holes};
  \draw[-{Stealth}, targetgreen!60!black] (eng.west) -- (3.15,0.35);
  \node[anchor=west, align=left, trapred] (wng) at (5.1,-1.3)
    {wings, fuselage, tail:\\ many holes};
  \draw[-{Stealth}, trapred] (wng.west) -- (4.2,-0.6);
\end{tikzpicture}
\caption{Hits recorded on bombers that returned (schematic). The data are conditioned on
survival. Areas that look safe on survivors are the ones where a hit was most likely to
bring the plane down, so the missing planes carry the important information.}
\label{fig:sampling-wald}
\end{figure}

Survivorship bias is everywhere in tech: studying only startups that succeeded, only
code that is still in the repository, only users who did not churn, only requests that
completed. Section~\ref{sec:sampling-logs} returns to the last one.

\subsection{Response bias}

\textbf{Response bias} is error in the answers rather than in who answers. Causes include
leading wording (``Don't you agree the new design is cleaner?''), social desirability
(people under-report drinking and over-report voting), recall error (``How many hours did
you use the app last month?''), and the order in which options appear. Logged behaviour
is usually more reliable than self-report for the same question. Wording effects can be
large: changing ``forbid'' to ``not allow'' in otherwise identical questions has shifted
survey answers by more than $20$ percentage points.

\section{Size does not cure bias: the Literary Digest poll}
\label{sec:sampling-digest}

In 1936 the magazine \emph{Literary Digest} mailed about $10$ million straw-poll ballots
(unofficial mock ballots asking ``whom will you vote for?'')
for the US presidential election between Franklin Roosevelt and Alf Landon. About
$2.4$ million came back. The Digest predicted Landon would win with $57\%$ to Roosevelt's
$43\%$. Roosevelt won with about $61\%$ of the vote, one of the largest landslides in US
history.

George Gallup, with a sample of about $50{,}000$ chosen to mirror the electorate,
predicted a Roosevelt win with $56\%$. He also predicted, in advance, that the Digest
would show Roosevelt near $44\%$.

\begin{figure}[H]
\centering
\begin{tikzpicture}
\begin{axis}[
  width=10cm, height=5.6cm, ybar, bar width=1.1cm,
  ymin=0, ymax=75, ylabel={Roosevelt share (\%)},
  symbolic x coords={Digest,Gallup,Actual},
  xticklabels={{Digest: $43\%$\\($n\approx 2.4$M)},{Gallup: $56\%$\\($n\approx 50$k)},{Actual: $61\%$\\(election)}},
  xtick=data, enlarge x limits=0.25, axis lines*=left,
  x tick label style={font=\small, align=center},
]
\addplot[fill=histfill, draw=distblue] coordinates {(Digest,43) (Gallup,56) (Actual,61)};
\draw[densityred, dashed, thick] (rel axis cs:0,0.6667) -- (rel axis cs:1,0.6667);
\node[densityred, anchor=south west, font=\small, xshift=4pt] at (rel axis cs:0,0.6667) {50\%};
\end{axis}
\end{tikzpicture}
\caption{The largest poll was the worst by far. The Digest was off by about $18$
percentage points and called the wrong winner; Gallup, with a sample $48$ times smaller,
was off by about $5$ points.}
\label{fig:sampling-digest}
\end{figure}

What went wrong? The Digest built its mailing list from its own subscribers, telephone
directories, and car registrations. In 1936, during the Depression, phones and cars were
markers of wealth, and wealthier voters leaned Republican (selection bias). Then only
about a quarter of recipients replied, and Landon supporters were more motivated to
reply (nonresponse bias). Both biases pushed the same way.

How large was the Digest's \emph{chance} error? For a proportion near $0.43$ and
$n = 2.4$ million, the standard error (Chapter 6) is
\[
  \sqrt{\frac{0.43 \times 0.57}{2{,}400{,}000}} \approx 0.00032,
\]
about $0.03$ percentage points. The chance error was negligible. The entire $18$-point
miss was bias.

\begin{recoverybox}
Pause and think: if the Digest had mailed $100$ million ballots to the same kinds of
lists, would its prediction have improved? (No. The SE would drop to about $0.005$
points, and the estimate would sit even more firmly near $43\%$.)
\end{recoverybox}

\begin{figure}[H]
\centering
\begin{tikzpicture}
\begin{axis}[
  width=12cm, height=5.2cm, xmode=log, axis lines=left,
  xmin=100, xmax=10000000, ymin=0, ymax=42,
  xlabel={sample size $n$ (log scale)}, ylabel={error (percentage points)},
  legend style={font=\small, at={(0.98,0.98)}, anchor=north east},
]
\addplot[name path=bu, draw=none, forget plot] coordinates {
 (100,28.000) (178,25.499) (316,23.623) (562,22.217) (1000,21.162) (1778,20.371)
 (3162,19.778) (5623,19.334) (10000,19.000) (17783,18.750) (31623,18.562)
 (56234,18.422) (100000,18.316) (177828,18.237) (316228,18.178) (562341,18.133)
 (1000000,18.100) (1778279,18.075) (3162278,18.056) (5623413,18.042) (10000000,18.032)};
\addplot[name path=bl, draw=none, forget plot] coordinates {
 (100,8.000) (178,10.501) (316,12.377) (562,13.783) (1000,14.838) (1778,15.629)
 (3162,16.222) (5623,16.666) (10000,17.000) (17783,17.250) (31623,17.438)
 (56234,17.578) (100000,17.684) (177828,17.763) (316228,17.822) (562341,17.867)
 (1000000,17.900) (1778279,17.925) (3162278,17.944) (5623413,17.958) (10000000,17.968)};
\addplot[trapred!25] fill between[of=bu and bl];
\addlegendentry{biased sample: $18 \pm 2\,\SE$}
\addplot[trapred, thick, domain=100:10000000] {18};
\addlegendentry{bias $= 18$ points}
\addplot[distblue, very thick] coordinates {
 (100,10.000) (178,7.499) (316,5.623) (562,4.217) (1000,3.162) (1778,2.371)
 (3162,1.778) (5623,1.334) (10000,1.000) (17783,0.750) (31623,0.562) (56234,0.422)
 (100000,0.316) (177828,0.237) (316228,0.178) (562341,0.133) (1000000,0.100)
 (1778279,0.075) (3162278,0.056) (5623413,0.042) (10000000,0.032)};
\addlegendentry{random sample: $2\,\SE$}
\end{axis}
\end{tikzpicture}
\caption{Typical error versus sample size for a proportion near $50\%$. A random
sample's error ($2\,\SE = 100/\sqrt n$ points) falls toward zero. A sample with $18$
points of bias converges, with growing confidence, to the wrong answer. A random sample
of $1000$ beats a biased sample of any size.}
\label{fig:sampling-size-vs-bias}
\end{figure}

\begin{happeningbox}
Increasing $n$ shrinks variance, never bias. A huge biased sample is worse than useless
because its tiny standard error makes the wrong answer look precise. Big data from logs
or scraped sources is often exactly this: millions of rows, collected by a process
nobody randomised.
\end{happeningbox}

\begin{connectionbox}
The total error of an estimate combines both pieces. The \textbf{mean squared error}
(the average squared distance between estimate and parameter over repeated samples)
splits as $\MSE = \Bias^2 + \Var$. Chapter 14 proves this decomposition; Chapter 15 meets it again
as the bias--variance trade-off in machine learning. The Digest had
$\MSE \approx 18^2 + 0.03^2$, all bias.
\end{connectionbox}

\section{Observational studies and experiments}
\label{sec:sampling-obs-exp}

Sampling decides whom you measure. A second design choice decides whether you can claim
cause and effect.

\begin{definition}[Observational study, experiment]
In an \textbf{observational study} the researcher records the treatment each unit
happens to receive (a feature used, a drug taken, a habit) without influencing it. In an
\textbf{experiment} the researcher \emph{assigns} the treatment.
\end{definition}

\subsection{Confounding}

Your analytics show that users of a new ``saved searches'' feature have $85\%$ 30-day
retention, compared with $45\%$ for non-users. Does the feature cause a $40$-point lift?

Split users by engagement level, measured \emph{before} the feature launched:

\begin{center}
\renewcommand{\arraystretch}{1.2}
\begin{tabular}{@{}lcccc@{}}
\toprule
 & \multicolumn{2}{c}{Uses feature} & \multicolumn{2}{c}{Does not use} \\
\cmidrule(lr){2-3}\cmidrule(lr){4-5}
Group & users & retained & users & retained \\
\midrule
Heavy users & $900$ & $810$ ($90\%$) & $100$ & $90$ ($90\%$) \\
Light users & $100$ & $40$ ($40\%$) & $900$ & $360$ ($40\%$) \\
\midrule
All & $1000$ & $850$ ($85\%$) & $1000$ & $450$ ($45\%$) \\
\bottomrule
\end{tabular}
\end{center}

Within each engagement group, feature users and non-users retain at exactly the same
rate. The $40$-point gap comes entirely from \emph{who} adopts the feature: $90\%$ of
adopters are heavy users, who retain well anyway.

\begin{figure}[H]
\centering
\begin{tikzpicture}
\begin{axis}[
  width=11cm, height=5.4cm, ybar, bar width=0.6cm,
  ymin=0, ymax=110, ylabel={30-day retention (\%)},
  symbolic x coords={All,Heavy,Light},
  xticklabels={{All users},{Heavy users},{Light users}},
  xtick=data, enlarge x limits=0.25, axis lines*=left,
  nodes near coords, nodes near coords style={font=\footnotesize},
  legend style={font=\small, at={(0.98,0.98)}, anchor=north east},
]
\addplot[fill=quotientorange!70, draw=quotientorange!70!black] coordinates {(All,85) (Heavy,90) (Light,40)};
\addlegendentry{uses feature}
\addplot[fill=histfill, draw=distblue] coordinates {(All,45) (Heavy,90) (Light,40)};
\addlegendentry{does not use}
\end{axis}
\end{tikzpicture}
\caption{A $40$-point gap overall, zero gap within each engagement level. Engagement
drives both feature adoption and retention, so comparing adopters with non-adopters
mixes the feature's effect with the engagement effect.}
\label{fig:sampling-confound-bars}
\end{figure}

\begin{definition}[Confounder]
A \textbf{confounder} (a hidden third factor that pushes on both sides of a
comparison) is a variable that influences both the treatment and the outcome,
creating an association between them that is not caused by the treatment.
\end{definition}

\begin{figure}[H]
\centering
\begin{tikzpicture}[font=\small, node distance=1.5cm and 3.2cm,
  box/.style={draw, thick, rounded corners, minimum width=2.8cm, minimum height=0.8cm, align=center}]
  \node[box, fill=structurepurple!8, draw=structurepurple] (c) {engagement\\(confounder)};
  \node[box, fill=quotientorange!8, draw=quotientorange, below left=of c] (t) {uses feature\\(treatment)};
  \node[box, fill=targetgreen!8, draw=targetgreen, below right=of c] (y) {retained\\(outcome)};
  \draw[-{Stealth}, thick, structurepurple] (c) -- (t);
  \draw[-{Stealth}, thick, structurepurple] (c) -- (y);
  \draw[-{Stealth}, thick, dashed, gray] (t) -- node[below]{causal effect? (here: none)} (y);
\end{tikzpicture}
\caption{The confounding triangle. Arrows mean ``influences''. The observed association
between treatment and outcome flows through the back path via engagement, whether or
not the direct arrow exists.}
\label{fig:sampling-confound-dag}
\end{figure}

\begin{recoverybox}
Pause and think: in the table, what fraction of heavy users adopted the feature, and what
fraction of light users? ($900/1000 = 90\%$ and $100/1000 = 10\%$.) That imbalance,
combined with heavy users retaining better, is the whole $40$-point gap.
\end{recoverybox}

Splitting by engagement fixed this example because we knew the confounder and had
measured it. In general you cannot list every confounder. Heavy users differ from light
users in a hundred ways: job, device, motivation, how they found the product. Adjusting
for the ones you thought of leaves the ones you did not. Chapter 13 shows how multiple
regression adjusts for measured confounders and why it cannot fix unmeasured ones.

\begin{trapbox}
``Correlation is not causation'' is too weak a slogan. The useful version: an
observational difference between groups equals the treatment effect \emph{plus} the
pre-existing difference between the kinds of units that chose each treatment. You see
only the sum.
\end{trapbox}

\subsection{Randomised experiments}

The fix is to take the choice away from the units. Flip a coin for each user: heads gets
the feature, tails does not. Now heavy and light users, motivated and unmotivated, iOS
and Android, all land in each group in roughly equal proportions. Any confounder, measured
or not, is balanced \emph{on average}.

\begin{figure}[H]
\centering
\begin{tikzpicture}[font=\small, node distance=0.9cm and 1.1cm,
  box/.style={draw, thick, rounded corners, minimum height=0.85cm, align=center, inner sep=4pt}]
  \node[box, fill=gray!8] (pool) {$2000$ eligible\\users};
  \node[box, fill=structurepurple!10, draw=structurepurple, right=of pool] (rand) {random\\assignment};
  \node[box, fill=quotientorange!10, draw=quotientorange, above right=0.35cm and 1.1cm of rand] (trt) {Treatment\\$n=1000$, feature on};
  \node[box, fill=histfill!25, draw=distblue, below right=0.35cm and 1.1cm of rand] (ctl) {Control\\$n=1000$, feature off};
  \node[box, fill=gray!8, right=4.3cm of rand] (cmp) {compare\\retention};
  \draw[-{Stealth}, thick] (pool) -- (rand);
  \draw[-{Stealth}, thick] (rand) -- (trt.west);
  \draw[-{Stealth}, thick] (rand) -- (ctl.west);
  \draw[-{Stealth}, thick] (trt.east) -- (cmp);
  \draw[-{Stealth}, thick] (ctl.east) -- (cmp);
\end{tikzpicture}
\caption{A randomised controlled experiment. Because a random mechanism, not the user,
decides who gets the feature, the two groups differ systematically only in the
treatment. A difference in outcomes beyond chance is therefore caused by it.}
\label{fig:sampling-randomise}
\end{figure}

\begin{whybox}
Before treatment, the two random groups are two random samples from the same pool. So
their average engagement, age, device mix, and every other pre-treatment trait differ
only by chance, and chance differences shrink as the groups grow. After treatment, if
outcomes differ by more than chance can explain (Chapter 9 makes ``more than chance''
precise), the only systematic difference left is the treatment itself.
\end{whybox}

The essential ingredients of a good experiment:
\begin{itemize}
  \item \textbf{Control group.} A comparison group that does not get the treatment.
  Without it, a before/after change could be seasonality, a marketing push, or a bug fix
  shipped the same week.
  \item \textbf{Randomisation.} Assignment by a random mechanism, never by judgement,
  convenience, or the units' own choice.
  \item \textbf{Replication.} Enough units that chance differences between groups are
  small relative to the effect you want to detect.
\end{itemize}

\subsection{Blocking}

Randomisation balances groups on average, but a particular assignment can be lopsided.
If iOS users convert at $8\%$ and Android users at $3\%$, a random split that happens to
put $55\%$ of iOS users in the treatment group adds noise. \textbf{Blocking} is
stratification applied to experiments: group units by a known, important trait, then
randomise separately within each block.

\begin{figure}[H]
\centering
\begin{tikzpicture}[font=\small, node distance=0.6cm and 0.9cm,
  box/.style={draw, thick, rounded corners, minimum height=0.75cm, align=center, inner sep=3pt}]
  \node[box, fill=gray!8] (all) {all users};
  \node[box, fill=sourceblue!10, draw=sourceblue, above right=0.9cm and 1.0cm of all] (ios) {block: iOS};
  \node[box, fill=sourceblue!10, draw=sourceblue, below right=0.9cm and 1.0cm of all] (and) {block: Android};
  \node[box, fill=structurepurple!10, draw=structurepurple, right=of ios] (r1) {randomise};
  \node[box, fill=structurepurple!10, draw=structurepurple, right=of and] (r2) {randomise};
  \node[box, fill=quotientorange!10, draw=quotientorange, above right=0.05cm and 0.9cm of r1] (t1) {A};
  \node[box, fill=histfill!25, draw=distblue, below right=0.05cm and 0.9cm of r1] (c1) {B};
  \node[box, fill=quotientorange!10, draw=quotientorange, above right=0.05cm and 0.9cm of r2] (t2) {A};
  \node[box, fill=histfill!25, draw=distblue, below right=0.05cm and 0.9cm of r2] (c2) {B};
  \node[box, fill=gray!8, right=3.2cm of r1, yshift=-1.25cm] (cmp) {combine\\within-block\\differences};
  \draw[-{Stealth}, thick] (all) -- (ios.west);
  \draw[-{Stealth}, thick] (all) -- (and.west);
  \draw[-{Stealth}, thick] (ios) -- (r1);
  \draw[-{Stealth}, thick] (and) -- (r2);
  \foreach \r/\a/\b in {r1/t1/c1, r2/t2/c2}{
    \draw[-{Stealth}, thick] (\r.east) -- (\a.west);
    \draw[-{Stealth}, thick] (\r.east) -- (\b.west);
    \draw[-{Stealth}, gray] (\a.east) -- (cmp.west);
    \draw[-{Stealth}, gray] (\b.east) -- (cmp.west);}
\end{tikzpicture}
\caption{A randomised block design. Platform is fixed before randomising, so A and B
contain exactly the same platform mix. Platform differences cancel out of the
comparison instead of inflating its noise.}
\label{fig:sampling-blocking}
\end{figure}

The slogan: \emph{block what you can, randomise what you cannot}. Blocking handles the
few big known factors; randomisation handles everything else.

\subsection{Matched pairs}

The extreme form of blocking is a block of size two. To test a new kernel setting, run
each of $6$ servers once with the old setting and once with the new, in random order.
Measured mean latencies (ms):

\begin{center}
\renewcommand{\arraystretch}{1.15}
\begin{tabular}{@{}lcccccc@{}}
\toprule
Server & 1 & 2 & 3 & 4 & 5 & 6 \\
\midrule
Old & 120 & 95 & 210 & 150 & 80 & 175 \\
New & 112 & 90 & 198 & 146 & 75 & 163 \\
Difference & 8 & 5 & 12 & 4 & 5 & 12 \\
\bottomrule
\end{tabular}
\end{center}

The old and new means are $138.3$ and $130.7$ ms; the SDs across servers are about $49$
and $47$ ms. Compared as two separate groups, a $7.7$ ms gap is buried in $48$ ms of
server-to-server spread. But every server improved. The differences have mean
$46/6 \approx 7.67$ ms and SD about $3.6$ ms. Pairing removes the server-to-server
variation from the comparison.

\begin{figure}[H]
\centering
\begin{tikzpicture}
\begin{axis}[
  width=7.5cm, height=6.2cm, axis lines=left,
  xmin=-0.3, xmax=1.3, ymin=60, ymax=225,
  xtick={0,1}, xticklabels={old setting, new setting},
  ylabel={mean latency (ms)},
]
\foreach \a/\b in {120/112, 95/90, 210/198, 150/146, 80/75, 175/163}{
  \edef\tmp{\noexpand\addplot[distblue, thick, mark=*, mark size=2pt] coordinates {(0,\a) (1,\b)};}
  \tmp}
\end{axis}
\end{tikzpicture}
\caption{Matched pairs: each line is one server. Levels differ hugely between servers
($75$ to $210$ ms), but every line slopes down by $4$ to $12$ ms. Analysing the
differences uses that consistency; comparing the two columns as unrelated groups throws
it away.}
\label{fig:sampling-pairs}
\end{figure}

Matched pairs also covers before/after measurements on the same unit (with a
randomised order or a control to rule out time trends), and twin studies. Chapter 9
gives the paired $t$-test that analyses such data.

\subsection{Blinding and placebos}

People who know they are treated behave differently. A patient told they are getting a
new painkiller may report less pain whether or not the pill does anything: the
\textbf{placebo effect}. A doctor who knows which patients got the drug may rate them
more favourably.

\begin{itemize}
  \item A \textbf{placebo} is a fake treatment (a sugar pill, a sham procedure) given to
  the control group so both groups have the same experience apart from the active
  ingredient.
  \item \textbf{Single-blind}: subjects do not know their group.
  \textbf{Double-blind}: neither subjects nor the people measuring outcomes know.
\end{itemize}

In software the analogue is usually automatic: users rarely know they are in an
experiment. Blinding matters for the \emph{analysts}. Deciding which metric is
``primary'', when to stop, and which users to exclude, \emph{after} seeing which arm is
winning, reintroduces judgement into a design built to exclude it (Chapter 11).

\section{A/B tests are randomised experiments}
\label{sec:sampling-ab}

An A/B test is the experiment of Figure~\ref{fig:sampling-randomise} at web scale. Each
user is assigned to variant A (control) or B (treatment) by a deterministic function of
their user ID that behaves like a coin flip. For example, A has $10{,}000$ users with
$520$ conversions ($5.2\%$) and B has $10{,}000$ users with $580$ ($5.8\%$). Because
assignment was random, the $0.6$-point difference is either caused by B or is chance;
Chapters 9 and 10 decide which.

\begin{lstlisting}[language=Python]
import hashlib
def variant(user_id: str, experiment: str = "checkout_v2") -> str:
    h = hashlib.sha256(f"{experiment}:{user_id}".encode()).hexdigest()
    return "B" if int(h, 16) % 100 < 50 else "A"   # 50/50 split
# Deterministic per user (same user, same variant on every visit),
# but unrelated to anything about the user: effectively random.
\end{lstlisting}

\begin{trapbox}
Common ways A/B tests stop being randomised experiments:
\begin{itemize}
  \item \textbf{Analysing only users who saw the change.} If B's new button is below the
  fold (off the first screen, visible only after scrolling), ``users who saw it'' are the users who scrolled, a self-selected group.
  Compare everyone \emph{assigned} to each arm, whether or not they saw the change.
  This is called \textbf{intention-to-treat} analysis: you analyse the groups the coin
  created, so the comparison stays randomised.
  \item \textbf{Sample ratio mismatch.} You configured $50/50$ but got $10{,}000$ vs
  $9{,}700$. Something (a crash, a redirect, a bot filter) removed users from one arm
  non-randomly. Investigate before reading any metric.
  \item \textbf{Comparing to last week instead of a concurrent control.} That is a
  before/after observational study.
\end{itemize}
\end{trapbox}

\section{Generalisation versus causation}
\label{sec:sampling-2x2}

The two kinds of randomness answer different questions:
\begin{itemize}
  \item \textbf{Random sampling} (who is in the study) licenses \emph{generalisation}
  from the sample to the population.
  \item \textbf{Random assignment} (who gets which treatment) licenses \emph{causal}
  conclusions about the treatment.
\end{itemize}
They are independent. A study can have both, either, or neither.

\begin{figure}[H]
\centering
\begin{tikzpicture}[font=\small]
  \def\w{4.6} \def\h{2.3}
  \fill[targetgreen!18] (0,\h) rectangle (\w,2*\h);
  \fill[sourceblue!12] (\w,\h) rectangle (2*\w,2*\h);
  \fill[quotientorange!14] (0,0) rectangle (\w,\h);
  \fill[gray!10] (\w,0) rectangle (2*\w,\h);
  \draw[thick] (0,0) rectangle (2*\w,2*\h);
  \draw[thick] (\w,0) -- (\w,2*\h); \draw[thick] (0,\h) -- (2*\w,\h);
  \node[align=center, text width=4.2cm] at (0.5*\w,1.5*\h)
    {\textbf{Causal + generalisable}\\A/B test on all traffic;\\trial on a random sample\\of customers};
  \node[align=center, text width=4.2cm] at (1.5*\w,1.5*\h)
    {\textbf{Generalisable, not causal}\\random survey of users;\\opinion poll};
  \node[align=center, text width=4.2cm] at (0.5*\w,0.5*\h)
    {\textbf{Causal for these units only}\\clinical trial on volunteers;\\A/B test on beta users};
  \node[align=center, text width=4.2cm] at (1.5*\w,0.5*\h)
    {\textbf{Neither}\\``users of feature X retain more''\\from whatever logs exist};
  \node[anchor=south] at (0.5*\w,2*\h+0.05) {random assignment: \textbf{yes}};
  \node[anchor=south] at (1.5*\w,2*\h+0.05) {random assignment: \textbf{no}};
  \node[anchor=east, align=right] at (-0.1,1.5*\h) {random\\sampling:\\ \textbf{yes}};
  \node[anchor=east, align=right] at (-0.1,0.5*\h) {random\\sampling:\\ \textbf{no}};
\end{tikzpicture}
\caption{What each design licenses. Rows: were units randomly sampled from the target
population? Columns: were treatments randomly assigned? Most medical trials sit
bottom-left (causal, but only for people like the volunteers); most dashboards sit
bottom-right.}
\label{fig:sampling-2x2}
\end{figure}

\begin{intuitionbox}
An A/B test randomises assignment, so it is causal. Whether it generalises depends on
who was in it. A test run during the holiday week, or only in one country, or only on
the beta channel, tells you the effect for that group. Extending it to everyone is an
assumption, not a conclusion.
\end{intuitionbox}

\begin{recoverybox}
Pause and think: a company emails a random $1\%$ of all customers a discount, and
compares their purchases with the other $99\%$. Which cell is it in? (Top-left. Every
customer is in the study, so the ``sample'' is the whole population, and the discount
was assigned at random. The result is causal and applies to all current customers.)
\end{recoverybox}

\section{Selection in tech logs and data leakage}
\label{sec:sampling-logs}

Logs feel like a census: every request, every user, nothing sampled. They are not. Logs
are produced by code, and code decides what gets written. That decision is a selection
mechanism, and it is rarely random.

\subsection{Survivorship in latency logs}

A service handles $1000$ requests. Clients give up at $2000$ ms; the server only writes a
latency record when it sends a response, so the $40$ requests that timed out leave no
record. The $960$ logged latencies look like this:

\begin{figure}[H]
\centering
\begin{tikzpicture}
\begin{axis}[
  width=12cm, height=5.6cm, axis lines=left,
  xmin=0, xmax=2200, ymin=0, ymax=470,
  xlabel={latency (ms)}, ylabel={requests},
  xtick={0,500,1000,1500,2000}, scaled x ticks=false,
  x tick label style={/pgf/number format/1000 sep={}},
]
\addplot[ybar interval, fill=histfill, draw=distblue] coordinates {
  (0,150) (100,380) (200,250) (300,100) (400,40) (500,20) (600,7) (700,5)
  (800,5) (900,3) (1000,3)};
\addplot[ybar interval, fill=trapred!15, draw=trapred, dashed, thick] coordinates {
  (2000,40) (2100,40)};
\draw[-{Stealth}, trapred] (axis cs:1750,200) -- (axis cs:2040,55);
\node[trapred, align=center, anchor=south, font=\small] at (axis cs:1700,200)
  {40 timeouts:\\never logged};
\draw[densityred, thick] (axis cs:215,0) -- (axis cs:215,440);
\node[densityred, anchor=west, font=\small] at (axis cs:230,440) {logged mean $215$ ms};
\end{axis}
\end{tikzpicture}
\caption{The logged latency histogram (blue) is conditioned on the request completing.
The slowest $4\%$ of requests are missing. Every summary computed from the log (mean,
p95, p99) is biased low, and the tail percentiles are the most wrong.}
\label{fig:sampling-logs}
\end{figure}

Using bin midpoints, the logged mean is
\[
  \frac{50\cdot 150 + 150\cdot 380 + 250\cdot 250 + \dots + 950\cdot 3}{960}
  = \frac{206{,}400}{960} = 215 \text{ ms}.
\]
Each timed-out request took at least $2000$ ms, so the true mean is at least
\[
  \frac{206{,}400 + 40 \cdot 2000}{1000} = \frac{286{,}400}{1000} \approx 286 \text{ ms}.
\]
The logged p99 is the $951$st of $960$ sorted values, which falls in the $700$--$800$ ms
bin. The true p99 is the $990$th of $1000$; since $40$ requests ($4\%$) timed out, it is
at least $2000$ ms. The dashboard says p99 is under a second; users experience the
opposite.

\begin{trapbox}
Before trusting any metric from logs, ask: \emph{which events fail to produce a log
line?} Crashes, timeouts, users who bounce before analytics loads, clients with ad
blockers, and requests dropped by a load balancer are typical. Log the attempt as well
as the outcome, so failures are counted.
\end{trapbox}

\subsection{Other selection traps in product data}

\begin{itemize}
  \item \textbf{Active-user filters.} ``Average session length among daily active
  users'' rises when a change drives casual users away, even though no one's sessions
  got longer. The denominator changed.
  \item \textbf{Opt-in telemetry.} Users who enable diagnostics skew toward power users
  and developers.
  \item \textbf{Feedback loops.} A recommender only logs clicks on items it chose to
  show. Training the next model on those logs teaches it about the old model's choices,
  not about user taste in general.
\end{itemize}

\subsection{Data leakage}

\textbf{Data leakage} is the machine-learning form of the same error: the training data
contain information that will not be available, or will look different, when the model
is used. Two common forms:
\begin{itemize}
  \item \textbf{Target leakage.} A churn model uses ``number of support tickets in the
  last 30 days'' computed over a window that ends \emph{after} the cancellation. The
  feature partly records the outcome. Accuracy looks excellent offline and collapses in
  production.
  \item \textbf{Selection leakage.} A fraud model is trained only on transactions that
  a previous rule system flagged for review, because only those have labels. The model
  learns the rule system's blind spots as ``safe''.
\end{itemize}

\begin{connectionbox}
Chapter 15 treats train/test splits and cross-validation. Every one of those procedures
assumes the test data are a random sample from the same process the model will face. The
sampling ideas of this chapter are exactly what that assumption means, and leakage is
what happens when it fails.
\end{connectionbox}

\begin{labbox}
Simulate nonresponse bias. Each of $100{,}000$ customers is satisfied with probability
$0.6$. Satisfied customers respond with probability $0.5$, unsatisfied ones with
probability $0.2$. Compute $\hat p$ among responders. Then raise $n$ to $10^7$: does
$\hat p$ approach $0.6$?

\textit{Answer.} Responders: $0.6\cdot 0.5 = 0.30$ satisfied, $0.4\cdot 0.2 = 0.08$
unsatisfied, so $\hat p \to 0.30/0.38 \approx 0.789$. More customers make $\hat p$
settle more tightly around $0.789$, not $0.6$.
\end{labbox}

\begin{lstlisting}[language=Python]
import numpy as np
rng = np.random.default_rng(0)
n = 100_000
satisfied = rng.random(n) < 0.6
responds = rng.random(n) < np.where(satisfied, 0.5, 0.2)
print(satisfied.mean(), satisfied[responds].mean())   # ~0.600, ~0.789
\end{lstlisting}

\section{Summary}
\label{sec:sampling-summary}

\begin{itemize}
  \item A \textbf{parameter} ($\mu,\sigma,p,\beta$) describes the population; a
  \textbf{statistic} ($\xbar,s,\hat p,b$) is computed from a sample and estimates it.
  \item A \textbf{simple random sample} gives every subset the same chance; under it
  $\xbar$ is unbiased for $\mu$. Random selection beats expert judgement because its
  errors are unsystematic and measurable.
  \item \textbf{Stratified} sampling (SRS within known groups) improves precision;
  \textbf{cluster} sampling (whole groups) cuts cost but loses precision;
  \textbf{systematic} sampling fails on periodic lists; \textbf{convenience} samples
  support no inference.
  \item \textbf{Bias} does not shrink with $n$. Main sources: selection, nonresponse,
  survivorship, response bias. The 1936 Literary Digest poll ($2.4$ million responses,
  $18$ points wrong) lost to Gallup's far smaller, better-designed sample.
  \item \textbf{Observational} comparisons are vulnerable to \textbf{confounding}.
  \textbf{Randomised experiments} with a control group balance all confounders on
  average, measured or not.
  \item \textbf{Blocking} and \textbf{matched pairs} remove known sources of variation;
  \textbf{blinding} and \textbf{placebos} remove expectation effects.
  \item Random \textbf{sampling} $\Rightarrow$ generalisation. Random
  \textbf{assignment} $\Rightarrow$ causation. An A/B test is causal for the users it
  included.
  \item Logs are selected data. Ask what fails to get logged; check for leakage before
  trusting an offline model metric.
\end{itemize}

Random sampling makes $\xbar$ right on average. Chapter 6 asks the next question: how far
from $\mu$ is one particular $\xbar$ likely to be? The answer, the standard error, is what
turns a random sample into a margin of error.

\section{Exercises}
\label{sec:sampling-exercises}

\begin{enumerate}
  \item Name the sampling design.
  (a) A school district picks $10$ of its $60$ schools at random and surveys every
  teacher in them.
  (b) A bank samples $1\%$ of accounts at random within each of its five regions.
  (c) A QA engineer checks every $200$th build artifact, starting from a random one of
  the first $200$.
  (d) A product manager posts a poll in the company chat.

  \item For each number, say whether it is a parameter or a statistic and give its
  symbol. (a) Of $400$ randomly chosen sessions, $7.5\%$ ended in a crash. (b) The mean
  age of all $12{,}340$ employees on the HR roster is $34.1$. (c) In a sample of $50$
  houses, the fitted price increase per extra bedroom is \$18{,}000.

  \item A survey is sent to $500$ users; $150$ respond and $120$ of them say ``yes''.
  Give the range of possible values for the proportion of all $500$ who would say
  ``yes''. What would it be if non-responders said ``yes'' at half the responders' rate?

  \item A city has three districts making up $50\%$, $30\%$, and $20\%$ of households.
  A stratified sample gives mean monthly electricity use of $400$, $600$, and $1000$ kWh
  respectively. Estimate the city-wide mean. Why would the unweighted average of the three
  stratum means be wrong?

  \item An observational study finds that people who carry lighters have higher lung
  cancer rates. Name the likely confounder and draw the confounding triangle. What
  experiment would settle whether lighters cause cancer, and why is it unnecessary?

  \item A service's client timeout is $5$ s, and timed-out requests are not logged. The
  logged mean latency is $300$ ms, and $2\%$ of all requests time out. Give a lower bound
  for the true mean latency.

  \item A metric is sampled once per week, every Monday at 09:00, for a year. What can
  go wrong? Propose a better design with the same number of measurements.

  \item Place each study in the $2\times 2$ grid of
  Figure~\ref{fig:sampling-2x2}. (a) $5000$ customers chosen at random from all
  customers are each randomly shown one of two pricing pages. (b) Employees who
  volunteered are randomly split between two keyboard layouts. (c) A random sample of
  households reports how much TV their children watch and their school grades.

  \item For the Literary Digest poll, compute the standard error of $\hat p$ with
  $n = 2.4$ million and $\hat p = 0.43$. Explain why that number says nothing about how
  accurate the poll was.

  \item You want to test whether a new code-review tool reduces review time. Ten teams
  volunteer. Teams differ greatly in typical review time. Describe a matched-pairs or
  blocked design and say what it gains over randomly giving the tool to $5$ teams.
\end{enumerate}

\subsection*{Answers}

\begin{enumerate}
  \item (a) Cluster. (b) Stratified. (c) Systematic. (d) Convenience (voluntary
  response).
  \item (a) Statistic, $\hat p = 0.075$. (b) Parameter, $\mu = 34.1$ (the whole
  population is measured). (c) Statistic, $b = 18{,}000$ dollars per bedroom.
  \item Minimum $120/500 = 24\%$; maximum $(120 + 350)/500 = 94\%$. Responders say yes at
  $80\%$; at half that, $40\%$ of $350$ is $140$, giving $(120+140)/500 = 52\%$.
  \item $0.5\cdot 400 + 0.3\cdot 600 + 0.2\cdot 1000 = 200 + 180 + 200 = 580$ kWh. The
  unweighted average $(400+600+1000)/3 \approx 667$ kWh over-weights the small,
  high-usage district.
  \item Smoking causes both carrying a lighter and lung cancer. The experiment would
  randomly assign people to carry lighters; randomisation would balance smoking across
  groups. It is unnecessary because splitting by smoking status already answers the
  question: among smokers, and separately among non-smokers, carrying a lighter is not
  associated with cancer.
  \item True mean $\ge 0.98 \cdot 300 + 0.02 \cdot 5000 = 294 + 100 = 394$ ms.
  \item Weekly and daily cycles mean Monday 09:00 may be systematically unusual (e.g.\
  a weekly peak), so the estimate is biased. Better: $52$ times chosen at random over the
  year, or one random time within each week (stratified by week).
  \item (a) Random sampling and random assignment: causal and generalisable to all
  customers. (b) Random assignment only: causal for volunteers. (c) Random sampling
  only: generalisable association, not causal.
  \item $\sqrt{0.43\cdot 0.57/2{,}400{,}000} \approx 0.00032$, i.e.\ $0.03$ percentage
  points. The SE measures chance error only; the poll's error was $18$ points of bias,
  which the SE cannot see.
  \item Pair teams with similar typical review times (5 pairs) and flip a coin within
  each pair for who gets the tool; or measure each team with and without the tool in a
  random order. Comparing within pairs removes the large team-to-team differences, so a
  real effect is visible with far fewer teams. Note the result generalises only to teams
  like the volunteers.
\end{enumerate}
